\documentclass[12pt]{article}
\usepackage[utf8]{inputenc}
\usepackage[T1]{fontenc}
\usepackage{hyperref}
\title{Kaili, Unde Language Resources for K-12 Curriculum Development}
\author{Compiled by Claude Sonnet 5 (claude-sonnet-5) for Andr\'as Kornai}
\date{}
\begin{document}
\maketitle

\section{Introduction}
Unde Kaili (unz; Glottocode unde1235) is an Austronesian language of Central Sulawesi, Indonesia, in the Kaili dialect chain alongside Ledo, Da'a and Baras. The ASJP database lists about 35,700 speakers. Very little specific to Unde was found. No dictionary, Bible translation or school programme for Unde was found in the search results.

\section{Primary Educational Resources}
None found for Unde. The searches turned up no primer, textbook or local-content (muatan lokal) programme for it.

\section{Linguistic and Reference Materials}
\begin{itemize}
\item ASJP wordlist for Kaili Unde, compiled by Viktoria Smirnova and citing Barr et al. 1979, \emph{Languages of Central Sulawesi}: \url{https://asjp.clld.org/languages/KAILI_UNDE}
\item Wikipedia, Kaili language (overview of the varieties): \url{https://en.wikipedia.org/wiki/Kaili_language}
\item Wikipedia, Ledo Kaili language (the largest Kaili variety, for comparison): \url{https://en.wikipedia.org/wiki/Ledo_Kaili_language}
\item Wikipedia, Kaili--Pamona languages: \url{https://en.wikipedia.org/wiki/Kaili%E2%80%93Pamona_languages}
\end{itemize}

\section{Religious and Community Resources}
None found.

\section{Digital and Audio Resources}
None found.

\section{Curriculum Development Context}
A 2025 study in the journal Stilistika on Kaili use in Tawaeli, Palu, reports declining use among younger speakers and no regional regulation governing the language's development and preservation. That study concerns Kaili in general, not Unde in particular. Any Unde curriculum would be built from scratch and would probably need to draw on materials for the closely related Ledo Kaili. Community involvement and fieldwork would come first.
\end{document}